\documentclass{situationist-preprint}

\title{The Justification of Revolution}
\msid{POL-23-a}
\msfamily{POL-23}
\msdiscipline{Political Philosophy}
\mstype{Principal manuscript}
\mskeywords{revolution; right of resistance; extreme injustice; proportionality; moral luck; uncertainty}

\begin{document}
\maketitle

\begin{abstract}
Is revolution ever justified? Kant denied that there is any right of revolution; Locke held that a people may resist a government that has dissolved itself by tyranny. We argue that revolution is justified under conditions of extreme injustice, but that its justification depends on outcomes that cannot be known in advance. Extreme injustice is necessary, because only it suspends the duty to obey and outweighs the grave costs revolution imposes on others; lesser injustice calls for legal means and civil disobedience. It is not sufficient, because revolution is justified only if the order it produces is substantially less unjust than the one it replaces, by a margin proportionate to its costs. Whether that condition holds depends on the choices of many agents over many years, and is not knowable when the revolution begins. We distinguish the justification of a revolution, which is fact-relative, from the reasonableness of taking part in it, which is evidence-relative, and argue that the outcome-dependence of the former is an instance of moral luck, not a defect of the account.
\end{abstract}

\section{The problem}

\citet[§§211--243]{Locke1689} held that a government which acts against the trust by which it holds power, invading the property and liberty of the people, dissolves itself, and that the people may then resist and establish a new legislative. \citet[Doctrine of Right, General Remark~A]{Kant1797} held, on the contrary, that there can be no right of resistance against the supreme legislative authority, since such a right would presuppose a public authority above the sovereign entitled to judge between sovereign and people, which is a contradiction. \citet{Burke1790} argued that destroying inherited institutions in the name of abstract rights would produce worse tyranny than it removed; \citet{Arendt1963} contrasted the American revolution, which founded a durable constitution, with the French, which ended in terror.

Contemporary work assimilates revolution to war. \citet{Buchanan2013} and \citet{Finlay2015} apply to revolution the principles of just war theory \citep{Walzer1977}: just cause, last resort, reasonable prospect of success and proportionality. We accept this structure, with one amendment. Just war theory states the conditions as things the agent can assess in advance. We argue that the most important of them, proportionality of outcome, cannot be assessed in advance for revolutions, and that this is not a reason to deny that revolutions can be justified but a feature of what their justification is.

\section{Framework}

\begin{definition}[Revolution]\label{d:revolution}
A \emph{revolution} is a collective attempt, by means outside the law and including the threat or use of force, to replace a government or the basic structure of a society.
\end{definition}

\begin{definition}[Extreme injustice]\label{d:extreme}
A political order is \emph{extremely unjust} if it systematically and gravely violates the basic rights of many of those subject to it, such as their rights to life, bodily integrity and liberty of conscience, and offers them no effective legal means of redress.
\end{definition}

\begin{definition}[Justification]\label{d:just}
A revolution is \emph{justified} if the reasons for undertaking it, given the facts including its outcome, outweigh the reasons against. Taking part is \emph{reasonable} for an agent if, given the evidence available to her, the expected reasons for it outweigh the expected reasons against.
\end{definition}

Definition~\ref{d:just} adopts Parfit's distinction between fact-relative and evidence-relative senses of \enquote{ought} \citep[ch.~7]{Parfit2011}. The question \enquote{was the revolution justified?} is fact-relative; the question \enquote{was she right to take part, given what she knew?} is evidence-relative.

\begin{principle}[Just cause]\label{p:cause}
A revolution is justified only if the order it opposes is extremely unjust.
\end{principle}

\begin{principle}[Proportionate improvement]\label{p:prop}
A revolution is justified only if the order it produces is substantially less unjust than the one it replaces, by a margin proportionate to the costs the revolution imposes; and a revolution against an extremely unjust order that meets this condition is justified.
\end{principle}

\begin{assumption}[Uncertainty]\label{a:uncertain}
When a revolution begins, whether the order it will produce meets the condition of Principle~\ref{p:prop} cannot be known.
\end{assumption}

Principle~\ref{p:cause} has two grounds. The first is the duty to obey. On POL-03 there is a weak, defeasible duty to obey just laws, grounded in fairness; it lapses for laws that are seriously unjust, but under an order that is unjust without being extremely so, much of the law remains just and binding, and the order provides legal means of reform. Where the duty to obey binds and legal means exist, a breach of law is permissible only under the conditions of ETH-17: grave harm, minimal and public breach, acceptance of the legal consequences. A revolution violates all three of the last conditions. The second ground is cost. Revolutions kill, destroy institutions, and expose the innocent to violence and disorder; only the systematic violation of basic rights is grave enough to outweigh such costs.

Principle~\ref{p:prop} states proportionality as a condition on outcomes. A revolution that overthrows an extremely unjust order and replaces it with one as unjust, or worse, has imposed its costs for nothing; one that replaces it with a substantially more just order may have imposed them for a good sufficient to outweigh them. The condition is not success: a victorious revolution that installs a new tyranny fails it. The second clause says that the two conditions are jointly sufficient.

Assumption~\ref{a:uncertain} is empirical. The outcome of a revolution depends on the conduct of the old regime, of foreign powers, of rival revolutionary factions, and of the population, over a period of years; and the historical record shows that revolutions with similar causes and aims have issued in constitutional government, in renewed tyranny, and in civil war. The overthrow of the Portuguese dictatorship in 1974 produced a constitutional democracy at modest cost; the revolutions of 1789 and 1917 passed through terror, and the second produced, on most assessments, an order no less unjust than the one it overthrew. Their participants could not have known at the outset which course events would take.

\section{The argument}

\begin{theorem}\label{t:yes}
Revolution can be justified: a revolution against an extremely unjust order whose outcome satisfies the condition of Principle~\ref{p:prop} is justified.
\end{theorem}

\begin{proof}
Immediate from the second clause of Principle~\ref{p:prop}. That the antecedent can be satisfied is shown by cases, such as that of 1974, in which an extremely unjust order was replaced by a substantially more just one at proportionate cost.
\end{proof}

\begin{theorem}\label{t:unknown}
Whether a revolution is justified cannot be known when it begins.
\end{theorem}

\begin{proof}
By Principles~\ref{p:cause} and~\ref{p:prop}, a revolution against an extremely unjust order is justified if and only if its outcome satisfies the condition of Principle~\ref{p:prop}. By Assumption~\ref{a:uncertain}, this cannot be known when the revolution begins. Hence whether the revolution is justified cannot be known then either. (Where the order is not extremely unjust, the revolution is not justified by Principle~\ref{p:cause}, which may be known in advance; the theorem concerns the cases where justification is in question.)
\end{proof}

\begin{corollary}\label{c:reasonable}
Taking part in a revolution may be reasonable although the revolution turns out not to be justified, and unreasonable although it turns out to be justified.
\end{corollary}

\begin{proof}
Reasonableness depends on the evidence available when the agent acts, justification on the outcome. By Theorem~\ref{t:unknown}, the outcome is not fixed by the evidence available at the outset, so the two can diverge in either direction.
\end{proof}

The structure is that of moral luck. \citet{Williams1981} argued that the justification of Gauguin's leaving his family to paint depended on his success, which he could not know in advance, and \citet{Nagel1979} classified such cases as luck in the way one's projects turn out. The revolutionary's case is graver, since the costs fall on others, but the structure is the same: the agent can act reasonably, but whether what she does is justified is settled by the world.

\section{Objections}

\begin{objection}[Kant]
There can be no right of revolution, since a right to resist the sovereign presupposes a judge between sovereign and people, which the sovereign is by definition.
\end{objection}

\begin{reply}
Kant's argument assumes that the sovereign constitutes a juridical condition, a condition in which right can be determined at all. An extremely unjust order, in the sense of Definition~\ref{d:extreme}, does not: it offers no effective means of redress for the violation of basic rights, and in relation to those whose rights it violates it is closer to a state of nature than to a civil condition. This is Locke's point that such a government has dissolved itself, and it is part of why, on POL-05, Locke provides the better account of legitimate government. The argument does not claim a legal right of revolution; Definition~\ref{d:just} concerns moral justification.
\end{reply}

\begin{objection}[Action-guidance]
A justification that cannot be known in advance cannot guide action. Justification should be assessed by what the agent could know.
\end{objection}

\begin{reply}
Action is guided by the evidence-relative notion, reasonableness, which Corollary~\ref{c:reasonable} keeps distinct. But the question whether a revolution was justified is not only a question about its agents. It is a question about whether the deaths and suffering it caused were warranted, a question asked by those who bore them, and its answer depends on what those costs bought. To assess the revolution only by what its leaders could know would be to say that a revolution which issued in a terror worse than the order it replaced was justified, provided its leaders had reasonable hopes. That is a judgment about the leaders, not about the revolution.
\end{reply}

\begin{objection}[Burke]
If the outcome of revolution cannot be known, prudence forbids it: one should not destroy existing institutions on a gamble.
\end{objection}

\begin{reply}
Inaction is also a gamble. Under an extremely unjust order, the costs of the status quo are grave, continuing and known, and the outcome of not revolting is no more certain than the outcome of revolting, since such orders also collapse, harden or are overthrown by others. Uncertainty bears on both options and does not by itself favor either. What it does support is a duty of revolutionaries to reduce it where they can: by building institutions that constrain the post-revolutionary order before power is seized, by limiting violence, and by keeping open the means of correction, measures that raise the probability that the condition of Principle~\ref{p:prop} will be met.
\end{reply}


\section{What the result does not establish}

The argument does not establish a legal right of revolution, a threshold of probability above which participation is reasonable, or the permissibility of particular means, such as violence against noncombatants, which falls under the ethics of war.

\section{Conclusion}

Revolution is justified when it opposes extreme injustice and produces an order substantially more just, by a margin proportionate to its costs. The first condition can be known in advance; the second cannot, because it depends on outcomes shaped by many agents over many years. Revolutionaries can act reasonably on the evidence, but whether their revolution is justified is settled only by what it brings about.

\begin{thebibliography}{99}\small
\bibitem[Arendt(1963)]{Arendt1963} Arendt, H. (1963). \emph{On Revolution}. Viking Press.
\bibitem[Buchanan(2013)]{Buchanan2013} Buchanan, A. (2013). The ethics of revolution and its implications for the ethics of intervention. \emph{Philosophy \& Public Affairs}, 41(4), 291--323.
\bibitem[Burke(1790)]{Burke1790} Burke, E. (1790). \emph{Reflections on the Revolution in France}. J. Dodsley.
\bibitem[Finlay(2015)]{Finlay2015} Finlay, C.~J. (2015). \emph{Terrorism and the Right to Resist: A Theory of Just Revolutionary War}. Cambridge University Press.
\bibitem[Kant(1797)]{Kant1797} Kant, I. (1797). \emph{Die Metaphysik der Sitten}. Trans. M. Gregor as \emph{The Metaphysics of Morals}, Cambridge University Press, 1996.
\bibitem[Locke(1689)]{Locke1689} Locke, J. (1689). \emph{Two Treatises of Government}. Cited by section of the Second Treatise.
\bibitem[Nagel(1979)]{Nagel1979} Nagel, T. (1979). Moral luck. In \emph{Mortal Questions}. Cambridge University Press.
\bibitem[Parfit(2011)]{Parfit2011} Parfit, D. (2011). \emph{On What Matters}. Vol.~1. Oxford University Press.
\bibitem[Walzer(1977)]{Walzer1977} Walzer, M. (1977). \emph{Just and Unjust Wars: A Moral Argument with Historical Illustrations}. Basic Books.
\bibitem[Williams(1981)]{Williams1981} Williams, B. (1981). Moral luck. In \emph{Moral Luck: Philosophical Papers 1973--1980}. Cambridge University Press.
\end{thebibliography}

\end{document}
